\par\Needspace{6\baselineskip}\medskip
\SourceLeaf{157}\pdfbookmark[2]{112.}{article-112}
\begin{center}112.\end{center}
Sint $x''$, $y''$, $z''$ coordinatae alicuius loci tertii, atque $u''$ eius argumentum latitudinis, $r''$ radius vector. Designabimus quantitates
\[
r'r''\sin(u''-u'),\qquad rr''\sin(u''-u),\qquad rr'\sin(u'-u),
\]
quae sunt areae duplae triangulorum inter radium vectorem secundum et tertium, primum et tertium, primum et secundum, resp. per $n$, $n'$, $n''$. Habebuntur itaque pro $x''$, $y''$, $z''$ expressiones iis similes, quas in art. praec. pro $x$, $y$, $z$ et $x'$, $y'$, $z'$ tradidimus, unde adiumento lemmatis I art. 78 facile deducuntur aequationes sequentes:
\[
\begin{aligned}
0&=nx-n'x'+n''x''\\
0&=ny-n'y'+n''y''\\
0&=nz-n'z'+n''z''.
\end{aligned}
\]
Sint iam longitudines geocentricae corporis coelestis tribus illis locis respondentes $\alpha$, $\alpha'$, $\alpha''$; latitudines geocentricae $\beta$, $\beta'$, $\beta''$; distantiae a terra ad eclipticam proiectae $\delta$, $\delta'$, $\delta''$; porro respondentes longitudines heliocentricae terrae $L$, $L'$, $L''$; latitudines $B$, $B'$, $B''$, quas non statuimus $=0$, ut liceat, tum parallaxis rationem habere, tum, si placet, pro ecliptica quodvis aliud planum adoptare; denique $D$, $D'$, $D''$ distantiae terrae a Sole ad eclipticam proiectae. Quodsi tunc $x$, $y$, $z$ per $L$, $B$, $D$, $\alpha$, $\beta$, $\delta$ exprimuntur, similiterque coordinatae ad locum secundum et tertium spectantes, aequationes praecedentes sequentem formam induunt:
\[
\begin{aligned}
1)\quad 0&=n(\delta\cos\alpha+D\cos L)-n'(\delta'\cos\alpha'+D'\cos L')\\
&\qquad+n''(\delta''\cos\alpha''+D''\cos L'')\\[.5ex]
2)\quad 0&=n(\delta\sin\alpha+D\sin L)-n'(\delta'\sin\alpha'+D'\sin L')\\
&\qquad+n''(\delta''\sin\alpha''+D''\sin L'')\\[.5ex]
3)\quad 0&=n(\delta\operatorname{tang}\beta+D\operatorname{tang}B)-n'(\delta'\operatorname{tang}\beta'+D'\operatorname{tang}B')\\
&\qquad+n''(\delta''\operatorname{tang}\beta''+D''\operatorname{tang}B'').
\end{aligned}
\]
Si hic $\alpha$, $\beta$, $D$, $L$, $B$, quantitatesque analogae pro duobus reliquis locis, tamquam cognitae spectantur, aequationesque per $n$, vel per $n'$, vel per $n''$ dividuntur, quinque incognitae remanent, e quibus itaque duas eliminare, sive per duas quascunque tres reliquas determinare licet. Hoc modo illae tres aequationes ad conclusiones plurimas gravissimas viam sternunt, e quibus quasdam imprimis insignes hic evolvemus.
